\begin{textsample}{POS Dim 6 – sc – Score 32.00 – sc/sc\_em\_59.txt}  \label{ex:f6_pos_003}
III \textbf{TRILHAS} DE APROFUNDAMENTO DA ÁREA CIÊNCIAS DA NATUREZA E SUAS \textbf{TECNOLOGIAS} 3.1 \textbf{INTRODUÇÃO} O Novo Ensino Médio busca fortalecer as relações entre os componentes curriculares, sendo essa também a proposta para a área de Ciências da Natureza e suas \textbf{tecnologias}.
A estrutura curricular contempla a formação geral básica, com carga horária total máxima de 1.800 horas, e os Itinerários \textbf{Formativos} ( \textbf{IF} ), com carga horária total mínima de 1.200 horas.
Estes últimos compõem a parte flexível do currículo e compreendem quatro componentes : a ) Projeto de Vida ; b ) \textbf{Trilhas} de aprofundamento ; c ) Componentes Curriculares Eletivos ; e d ) Segunda Língua Estrangeira.
As Diretrizes Curriculares Nacionais para o ensino médio determinam, em seu artigo 12, que os itinerários \textbf{formativos} da área de Ciências da Natureza e suas \textbf{tecnologias} sejam organizados considerando : > [... ] o aprofundamento de conhecimentos \textbf{estruturantes} para aplicação de diferentes conceitos em contextos sociais e de trabalho, organizando arranjos curriculares que permitam estudos em astronomia, metrologia, física geral, clássica, molecular, quântica e mecânica, instrumentação, ótica, acústica, química dos produtos naturais, análise de fenômenos físicos e químicos, meteorologia e climatologia, microbiologia, imunologia e parasitologia, ecologia, nutrição, zoologia, dentre outros, considerando o contexto local e as possibilidades de oferta pelos sistemas de ensino ( BRASIL, 2018, p. 7 ).
Os itinerários \textbf{formativos} das áreas do conhecimento se articulam com a formação geral básica, contemplando quatro eixos \textbf{estruturantes} : investigação científica, processos criativos, \textbf{mediação} e intervenção sociocultural e \textbf{empreendedorismo}, “ a fim de garantir que os(as) estudantes experimentem diferentes situações de aprendizagem e desenvolvam um conjunto diversificado de habilidades relevantes para sua formação integral ” ( BRASIL, 2018, p. 5 ).
Nesta seção, serão abordados apenas aspectos que dizem respeito ao desenvolvimento das \textbf{trilhas} de aprofundamento.
A construção dessas \textbf{trilhas} considera os eixos \textbf{estruturantes} mencionados e retoma os pressupostos epistemológicos da Teoria Histórico-Cultural da PCSC ( SANTA CATARINA, 2014 ), conforme representado na \textbf{figura} 2. \textbf{Figura} 2 : Eixos \textbf{estruturantes} Fonte : Elaboração dos autores ( 2020 ).
Na área de Ciências da Natureza e suas \textbf{tecnologias}, temos quatro \textbf{trilhas} de aprofundamento construídas a partir da discussão entre professores elaboradores da rede estadual de ensino e colaboradores desta área de conhecimento, com vistas a oportunizar aos(às) estudantes o aprofundamento dos seus conhecimentos nesta área.
As \textbf{trilhas} de aprofundamento buscam expandir os aprendizados promovidos pela formação geral, em articulação com temáticas contemporâneas, sintonizadas com o contexto e os interesses dos(as) estudantes.
Além de fomentar o interesse pelas vocações científico-tecnológicas, os aprofundamentos permitem que os jovens já concluam o ensino médio com algum diferencial em sua formação para a própria inserção no mundo do trabalho.
Foram estruturadas em percursos com começo, meio e fim ; seu fluxo compreende os quatro eixos \textbf{estruturantes} ( \textbf{Figura} 2 ) e permite aos(às) estudantes se desenvolverem de forma integral, orgânica e progressiva, lidando com desafios cada vez mais complexos.
Também é interessante que cada etapa dessa jornada integre e articule os conhecimentos, habilidades, atitudes e valores adquiridos nas etapas anteriores, ou concomitantes.
O aprofundamento da área de Ciências da Natureza e suas \textbf{tecnologias} no ensino médio é indispensável para o enfrentamento de questões \textbf{emergentes} da sociedade contemporânea, que exige a compreensão de diferentes conceitos de Ciências da Natureza, \textbf{tecnologias}, ambiente e \textbf{sustentabilidade} necessários para desvelar as realidades e constituir processos críticos e reflexivos.
O ato de viver nessa sociedade exige cidadãos reflexivos, críticos, criativos, questionadores, curiosos, com espírito de investigação e protagonismo, capazes de tomar decisões em âmbito individual e coletivo com base em \textbf{argumentos} científicos.
Compreender conceitos de Biologia, Física e Química, que regem o funcionamento da natureza, pode auxiliar os(as) estudantes ( jovens, adultos, idosos ) a tomar decisões, por exemplo, sobre sua própria alimentação ou sobre a utilização e consumo de produtos tecnológicos.
As \textbf{trilhas} foram propostas a partir de temas abrangentes que aprofundam os conhecimentos dos três componentes curriculares, porém, de forma integrada e com o propósito de desenvolver diversas habilidades.
Deste modo, busca -se aprofundar conteúdos escolares, relacionando -os com situações do cotidiano e com problemáticas da vida dos(as) estudantes.
Nas \textbf{trilhas} apresentadas à área, utilizam -se palavras-chave para as habilidades específicas dos eixos \textbf{estruturantes}, objetivando uma leitura dinâmica e fluida com sugestões de objetos do conhecimento.
O quadro 11 apresenta a relação entre as habilidades associadas às competências gerais da BNCC e às habilidades específicas de cada eixo \textbf{estruturante}, que fazem parte das \textbf{trilhas} de aprofundamento da área de Ciências da Natureza e suas \textbf{tecnologias}.
Quadro 25 - Eixos \textbf{estruturantes} e habilidades das \textbf{trilhas} de aprofundamento associadas às palavras-chave : competências gerais da BNCC Investigação científica : pensar e fazer científico Habilidades dos itinerários \textbf{formativos} associadas às competências gerais da BNCC - Identificar, selecionar, \textbf{processar} e analisar dados, fatos e evidências com curiosidade, atenção, criticidade e ética, utilizando, inclusive, o apoio de \textbf{tecnologias} \textbf{digitais}. - Posicionar -se com base em critérios científicos, éticos e estéticos, utilizando dados, fatos e evidências para \textbf{respaldar} conclusões, opiniões e \textbf{argumentos} por meio de afirmações claras, \textbf{ordenadas}, coerentes e \textbf{compreensíveis}, sempre respeitando valores universais, como liberdade, democracia, justiça social, pluralidade, solidariedade e \textbf{sustentabilidade}. - Utilizar informações, conhecimentos e ideias \textbf{resultantes} de investigações científicas para criar ou propor soluções para \textbf{problemas} diversos.
Palavras-chave/Habilidades específicas dos itinerários \textbf{formativos} associadas aos eixos \textbf{estruturantes} da área de Ciências da Natureza e suas \textbf{tecnologias} - Investigar e analisar \textbf{problemas} — Investigar e analisar situações-problema e variáveis que interferem na dinâmica de fenômenos da natureza e/ou de processos tecnológicos, considerando dados e informações disponíveis em diferentes mídias, com ou sem o uso de dispositivos e \textbf{aplicativos} \textbf{digitais}. - Levantar e \textbf{testar} hipóteses — Levantar e \textbf{testar} hipóteses sobre variáveis que interferem na dinâmica de fenômenos da natureza e/ou de processos tecnológicos, com ou sem o uso de dispositivos e \textbf{aplicativos} \textbf{digitais}, utilizando procedimentos e linguagens adequados à investigação científica. - Coletar e tratar dados — Selecionar e sistematizar, com base em estudos e/ou pesquisas ( bibliográfica, exploratória, de campo, experimental, etc. ) em fontes confiáveis, informações sobre a dinâmica dos fenômenos da natureza e/ou de processos tecnológicos, identificando os diversos pontos de vista, posicionando -se mediante argumentação, com o cuidado de citar as fontes dos recursos utilizados na pesquisa e buscando apresentar conclusões com o uso de diferentes mídias.
Processos criativos : pensar e fazer criativo Habilidades dos itinerários \textbf{formativos} associadas às competências gerais da BNCC - Reconhecer e analisar diferentes manifestações criativas, artísticas e culturais, por meio de vivências presenciais e virtuais que ampliem a visão de mundo, a sensibilidade, a criticidade e a criatividade. - Questionar, modificar e adaptar ideias existentes e criar propostas, obras ou soluções criativas, originais ou inovadoras, avaliando e assumindo riscos para lidar com as incertezas e colocá -las em prática. - Difundir novas ideias, propostas, obras ou soluções por meio de diferentes linguagens, mídias e plataformas, analógicas e \textbf{digitais}, com confiança e \textbf{coragem}, assegurando que alcancem os interlocutores pretendidos.
Palavras-chave/Habilidades específicas dos itinerários \textbf{formativos} associadas aos eixos \textbf{estruturantes} da área de Ciências da Natureza e suas \textbf{tecnologias} - Refletir criativa/criticamente — Reconhecer produtos e/ou processos criativos por meio de fruição, vivências e reflexão crítica sobre a dinâmica dos fenômenos naturais e/ou de processos tecnológicos, com ou sem o uso de dispositivos e \textbf{aplicativos} \textbf{digitais} ( como softwares de simulação e de realidade virtual, entre outros ). - Selecionar e mobilizar conhecimentos científicos — Selecionar e mobilizar intencionalmente recursos criativos relacionados às Ciências da Natureza para resolver \textbf{problemas} reais do ambiente e da sociedade, explorando e contrapondo diversas fontes de informação. - Propor e \textbf{testar} soluções — Propor e \textbf{testar} soluções éticas, estéticas, criativas e inovadoras para \textbf{problemas} reais, considerando a aplicação de design de soluções e o uso de \textbf{tecnologias} \textbf{digitais}, programação e/ou pensamento computacional que apoiem a construção de protótipos, dispositivos e/ou equipamentos, com o intuito de melhorar a qualidade de vida e/ou os processos produtivos. \textbf{Mediação} e intervenção sociocultural : convivência e atuação sociocultural Habilidades dos itinerários \textbf{formativos} associadas às competências gerais da BNCC - Reconhecer e analisar questões sociais, culturais e \textbf{ambientais} diversas, identificando e incorporando valores importantes para si e para o coletivo, que assegurem a tomada de decisões conscientes, \textbf{consequentes}, colaborativas e responsáveis. - Compreender e considerar a situação, a opinião e o sentimento do outro, agindo com empatia, flexibilidade e resiliência para promover o diálogo, a colaboração, a \textbf{mediação} e a resolução de conflitos, o combate ao preconceito e a valorização da diversidade. - Participar ativamente da proposição, implementação e avaliação de soluções para \textbf{problemas} socioculturais e/ou \textbf{ambientais} em nível local, regional, nacional e/ou global, corresponsabilizando -se pela realização de ações e projetos voltados ao bem comum.
Palavras-chave/Habilidades específicas dos itinerários \textbf{formativos} associadas aos eixos \textbf{estruturantes} da área de Ciências da Natureza e suas \textbf{tecnologias} - Identificar e explicar fenômenos — Identificar e explicar questões socioculturais e \textbf{ambientais} relacionadas a fenômenos físicos, químicos e/ou biológicos. - Selecionar e mobilizar conhecimentos científicos — Selecionar e mobilizar intencionalmente conhecimentos e recursos das Ciências da Natureza para propor ações individuais e/ou coletivas de \textbf{mediação} e intervenção sobre \textbf{problemas} socioculturais e \textbf{problemas} \textbf{ambientais}. - Propor e \textbf{testar} intervenções socioculturais e \textbf{ambientais} — Propor e \textbf{testar} estratégias de \textbf{mediação} e intervenção para resolver \textbf{problemas} de natureza sociocultural e de natureza \textbf{ambiental} relacionados às Ciências da Natureza. \textbf{Empreendedorismo} : autoconhecimento, \textbf{empreendedorismo} e projeto de vida Habilidades dos itinerários \textbf{formativos} associadas às competências gerais da BNCC - Reconhecer e utilizar qualidades e fragilidades pessoais com confiança para superar desafios e alcançar objetivos pessoais e profissionais, agindo de forma proativa e empreendedora, \textbf{perseverando} em situações de \textbf{estresse}, frustração, fracasso e \textbf{adversidade}. - Utilizar estratégias de planejamento, organização e \textbf{empreendedorismo} para estabelecer e adaptar metas, identificar caminhos, mobilizar apoios e recursos para realizar projetos pessoais e produtivos com foco, \textbf{persistência} e efetividade. - Refletir continuamente sobre seu próprio desenvolvimento e sobre seus objetivos presentes e futuros, identificando aspirações e oportunidades, inclusive relacionadas ao mundo do trabalho, que orientem escolhas, esforços e ações em relação à sua vida pessoal, profissional e cidadã.
Palavras-chave/Habilidades específicas dos itinerários \textbf{formativos} associadas aos eixos \textbf{estruturantes} da área de Ciências da Natureza e suas \textbf{tecnologias} - Avaliar possíveis soluções — Avaliar como oportunidades, conhecimentos e recursos relacionados às Ciências da Natureza podem ser utilizados na concretização de projetos pessoais ou produtivos, considerando as diversas \textbf{tecnologias} disponíveis e os impactos socioambientais. - Desenvolver projetos — Selecionar e mobilizar intencionalmente conhecimentos e recursos das Ciências da Natureza para desenvolver um projeto pessoal ou um empreendimento produtivo. - Criar propostas articuladas ao projeto de vida — Desenvolver projetos pessoais ou produtivos, utilizando as Ciências da Natureza e suas \textbf{tecnologias} para formular propostas concretas, articuladas com o projeto de vida.
Fonte : BRASIL ( 2018 ).
Ao final do ensino médio, espera -se que os(as) estudantes sejam capazes de realizar pesquisas científicas, criar obras, soluções e/ou inovações de modo a intervir positivamente na realidade, provocando, assim, mudanças significativas em sua vida e na comunidade em que estão inseridos.
Mobilizar, articular, integrar, calcular, investigar, comunicar, coletar, analisar, exercitar, ler, escrever, produzir são exemplos de verbos relacionados às habilidades que se esperam para a área.
Neste sentido, busca -se superar um ensino de Ciências centrado no resultado científico e na memorização de uma lista de conteúdos vinculados a séries específicas¹ que resultam em baixa participação dos(as) estudantes, em passividade, pouco interesse em aprender Ciências da Natureza, e na falta da devida preparação dos(as) jovens para sobreviver na sociedade contemporânea.
Alinha -se com uma educação atual, que pressupõe um(uma) estudante capaz de autogerenciar ou autogovernar seu processo de formação, fornecendo subsídios para que, ao longo da sua vida, seja capaz de buscar a autonomia ( MITRE et al., 2008 ).
Salienta -se que as orientações metodológicas apresentadas nas \textbf{trilhas} são coerentes com o ensino \textbf{desenvolvimental} de Davydov, que centraliza o ensino na teoria da atividade, pressuposto da Teoria Histórico-Cultural.
Recomenda -se consulta ao texto da Formação Geral Básica ( Ciências da Natureza e suas \textbf{tecnologias} ), que apresenta orientações relativas à avaliação e às metodologias ( experimentação, abordagem CTSA, ensino por investigação, aprendizagem baseada em projetos, mapas conceituais, problematizações, resolução de \textbf{problemas}, ensino por investigação, entre outras ).
As \textbf{trilhas} que apresentamos para a área — a ) Diálogos com nossas cidades : meio ambiente e \textbf{sustentabilidade} ; b ) Ciência, sociedade, convívio e saúde ; c ) A \textbf{tecnologia} das coisas : uma perspectiva sustentável na sociedade contemporânea ; e d ) Eu, nós e nossas escolhas : diálogos com a ciência para a transição das sociedades sustentáveis — foram produzidas para atender às demandas apresentadas pelos professores da área de Ciências da Natureza das escolas.
Cada uma das \textbf{trilhas} foi organizada em unidades curriculares.
As \textbf{trilhas} são semestrais, com carga horária de 160 horas, podendo ser ampliadas até 240 horas, dependendo da \textbf{matriz} curricular adotada pela unidade escolar e prevista no projeto político-pedagógico da instituição educativa.
A \textbf{trilha} “ Diálogos com nossas cidades - meio ambiente e \textbf{sustentabilidade} ” apresenta aos(às) estudantes a possibilidade de aprofundar conhecimentos da área de Ciências da Natureza, a partir da discussão da dinâmica das cidades.
Propõe a investigação de \textbf{problemas} ligados a planejamento urbano, eletricidade, poluição, mobilidade urbana, utilização consciente dos recursos naturais, entre outros.
Ao passar pela \textbf{trilha}, espera -se que os(as) estudantes conheçam as cidades e proponham alternativas de solução para \textbf{problemas} locais e promovam a interferência direta na comunidade.
A \textbf{trilha} “ Eureka!
Investigação no mundo da ciência ” tem como foco aprofundar conceitos da saúde associados à área de Ciências da Natureza e suas \textbf{tecnologias}, com abordagem de quatro unidades curriculares, envolvendo contexto local, regional e mundial por meio do processo de produção e divulgação dos conhecimentos científicos, do resgate dos conhecimentos tradicionais dos povos e suas contribuições para a ciência na área da saúde, de padrões “ estéticos ” “ supervalorizados ” pela sociedade atual, que causam implicações ( tanto externa quanto internamente ), relacionadas às vulnerabilidades e ao mau uso de \textbf{tecnologias} cotidianas, além de políticas de saúde e novas \textbf{tecnologias} que contribuem para a saúde e o bem-estar dos povos e geram novas profissões.
A \textbf{trilha} “ A \textbf{Tecnologia} das Coisas : uma perspectiva sustentável na sociedade contemporânea ” tem como foco a educação científica dos(as) estudantes e seu objetivo visa à problematização das diferentes concepções de \textbf{tecnologia} e ao incentivo à cultura \textbf{digital}, garantindo -lhes acessibilidade a novos conhecimentos mediados pela \textbf{tecnologia}, para a proposição de soluções criativas exigidas por seu contexto de vida.
Além disso, reflete sobre seus meios de produção e sobre a inserção de \textbf{tecnologias} energéticas que garantam a \textbf{sustentabilidade} do Planeta, mobilizando os(as) estudantes a avaliar seus impactos sociais, culturais e \textbf{ambientais}.
Na \textbf{trilha} “ Eu, nós e nossas escolhas : diálogos com a ciência para a transição das Sociedades Sustentáveis ”, a ênfase está na educação \textbf{ambiental} como um dos princípios norteadores, utilizando -se da metáfora das “ 5 CASAS ” ( TRAJBER, 2011 ), tendo por objetivo geral promover a educação para o processo de transição para uma sociedade sustentável, promovendo a compreensão de que o cuidado com o Planeta perpasse primeiramente pelo cuidado comigo ( EU - CORPO ), com o outro ( NÓS - FAMÍLIA ) e com o espaço em que vivemos ( COMUNIDADE - ESPAÇO SOCIAL ; TERRITÓRIO - ECOSSISTEMA ; PLANETA - NOSSA CASA COMUM ), numa perspectiva de interligação, de visão sistêmica da \textbf{sustentabilidade}.
Esse processo pode ser desenvolvido em espaços educativos sustentáveis. -- - ¹ O currículo do ensino médio do estado de Santa Catarina não distribui os objetos do conhecimento de cada componente curricular por série ( 1º, 2º e 3º ).
Os objetos são definidos para a área, tanto na formação geral básica, quanto nos itinerários \textbf{formativos} ( \textbf{trilhas} de aprofundamento, eletivas e projetos de vida ).
Esta flexibilidade do Novo Ensino Médio concede liberdade às escolas para definições da distribuição das habilidades e sugestões de objetos do conhecimento associados nos anos do ensino médio, em seus projetos político-pedagógicos.
A liberdade aos(às) estudantes está nas escolhas das \textbf{trilhas} de aprofundamento e nas disciplinas eletivas, a partir do que a escola oferece.

% matched lemmas: adversidade, ambiental, aplicativo, argumento, compreensível, consequente, coragem, desenvolvimental, digital, emergente, empreendedorismo, estresse, estruturante, figura, figurar, formativo, if, introdução, matriz, mediação, ordenado, perseverar, persistência, problema, processar, respaldar, resultante, sustentabilidade, tecnologia, testar, trilha
\end{textsample}
